\documentclass[10pt,a4paper]{amsart}
\usepackage[margin=19mm]{geometry}
\usepackage{amsmath,amssymb,mathtools}
\usepackage[scr=rsfs,cal=euler]{mathalfa}
\usepackage{xfrac}
\usepackage[unicode,hidelinks,bookmarks=false]{hyperref}
\newcommand{\ie}{i.e.\ }
\newcommand{\cf}{cf.\ }
\newcommand{\resp}{respectively\ }
\newcommand{\Subsection}{\subsection}
\providecommand{\nobreakdash}{\nobreak\hbox{-}}
\setlength{\emergencystretch}{2em}
\multlinegap0pt
\usepackage{amssymb}
\usepackage{mathtools}
\usepackage[shortlabels]{enumitem}
\newtheorem{theorem}{Theorem}
\newcommand{\R}{\mathbb R}
\newcommand{\norm}[1]{\left\|#1\right\|}
\title{Polylogarithmic Full-Chord Buffon Discrepancy}
\author{Samuel Korsky}
\date{Source: arXiv:2605.23020v1 (CC BY 4.0)}
\begin{document}
\maketitle

\noindent\emph{Source and licence.} Polylogarithmic Full-Chord Buffon Discrepancy. Version arXiv:2605.23020v1 (CC BY 4.0), distributed by arXiv under the Creative Commons Attribution 4.0 International licence, http://creativecommons.org/licenses/by/4.0/, as recorded in the arXiv licence metadata for this exact version and shown on the abstract page https://arxiv.org/abs/2605.23020v1. Full licence notice: \texttt{licenses/arxiv-2605.23020-CC-BY-4.0.txt}.\par\smallskip

\noindent\textit{Editorial scope.} Counted unit: the article's theorem thm:weighted-schmidt (source0.tex lines 136--152). The article's complete original proof of it is delivered with it (source0.tex lines 154--234, one \texttt{proof} environment of 81 lines). No mathematics is written, repaired or re-derived: the statement and the proof are the article's own lines, in the article's own notation, and no retained line is substituted or re-flowed.  Retained uncounted as the source-local context the proof needs: lines 124--134.  Nothing was omitted from any retained span, so the package carries no omission record.  No retained line contains a citation, so the wrapper carries no bibliography and no bibliography entry.  No retained line contains a figure environment, an image-inclusion command or drawing code, so no image is delivered and the package declares no asset dependency.  Editorial formatting only, in addition: an AMS \texttt{amsart} preamble that restates the article's own declarations and macros used by the retained lines, namely the environments \texttt{theorem} on separate counters, together with the standard packages those lines need.  The retained environments print in this document as Theorem 1 in order of appearance, whereas the article prints the same items with the numbering of its own full document.  The complete native source of the article as distributed in the arXiv e-print for this exact version is bundled unchanged as \texttt{sources/201177/source0.tex}.
\begin{theorem}[Schmidt rectangle lower bound]\label{thm:schmidt-rect}
There is an absolute constant $c_S>0$ such that for every axis-parallel square $Q\subset\R^2$ and every finite point set $P\subset Q$ with $N=\#P\ge2$,
\[
        \sup_{R\subset Q}
        \left|
        \#(P\cap R)-N\cdot\frac{|R|}{|Q|}
        \right|
        \ge c_S\log N,
\]
where the supremum is over all axis-parallel rectangles $R\subset Q$.
\end{theorem}

\begin{theorem}[Weighted Schmidt rectangle lower bound]\label{thm:weighted-schmidt}
Let $Q\subset\R^2$ be a fixed rectangle and let $\rho\in C^1(Q)$ satisfy
\[
        0<c_0\le \rho\le C_0<\infty,
        \qquad
        \norm{\rho}_{C^1(Q)}\le C_0.
\]
Then there are constants $c,M_0>0$, depending only on $Q,c_0,C_0$, with the following property. For every $M\ge M_0$ and every finite point set $P\subset Q$,
\[
        \sup_{R\subset Q}
        \left|
        \#(P\cap R)-M\int_R\rho(x,y)\,dx\,dy
        \right|
        \ge c\log M,
\]
where the supremum is over all axis-parallel rectangles $R\subset Q$.
\end{theorem}

\begin{proof}
Let
\[
        \Delta:=\sup_{R\subset Q}
        \left|
        \#(P\cap R)-M\int_R\rho
        \right|.
\]
We prove $\Delta\ge c\log M$ for large $M$.

\bigskip
\noindent
Choose a point $x_0$ in the interior of $Q$.  For $M$ large, let $Q_M\subset Q$ be the axis-parallel square centered at $x_0$ with side length
\[
        r=A\left(\frac{\log M}{M}\right)^{1/3},
\]
where $A>0$ will be chosen sufficiently small depending only on $c_S$ and $C_0$.  Put
\[
        N_M:=\#(P\cap Q_M).
\]
Since $Q_M$ is a rectangle,
\begin{equation}\label{eq:local-mass-control-schmidt}
        \left|N_M-M\int_{Q_M}\rho\right|\le \Delta.
\end{equation}
If $\Delta\ge (c_S/16)\log M$, we are done.  Hence assume the opposite.  Since $\rho\ge c_0$,
\[
        M\int_{Q_M}\rho
        \ge c_0Mr^2
        =c_0A^2M^{1/3}(\log M)^{2/3}.
\]
The assumed bound on $\Delta$ is negligible compared to this quantity for large $M$.  Therefore
\[
        N_M\asymp Mr^2,
        \qquad
        \log N_M\ge \frac14\cdot\log M
\]
for all sufficiently large $M$.

\bigskip
\noindent
Apply Theorem~\ref{thm:schmidt-rect} inside the square $Q_M$ to the point set $P\cap Q_M$.  There exists a rectangle $R\subset Q_M$ such that
\begin{equation}\label{eq:local-schmidt}
        \left|
        \#(P\cap R)-N_M\cdot\frac{|R|}{|Q_M|}
        \right|
        \ge \frac{c_S}{4}\cdot\log M.
\end{equation}
We now compare the constant-density target $N_M|R|/|Q_M|$ to the weighted target $M\int_R\rho$.  First, by \eqref{eq:local-mass-control-schmidt},
\[
        \left|
        \frac{|R|}{|Q_M|}\cdot N_M
        -
        \frac{|R|}{|Q_M|}\cdot M\int_{Q_M}\rho
        \right|
        \le \Delta.
\]
Second, if
\[
        \bar\rho_{Q_M}:=\frac1{|Q_M|}\int_{Q_M}\rho,
\]
then the $C^1$ bound on $\rho$ implies that $\rho$ varies by at most $O(C_0r)$ on $Q_M$. Therefore
\[
\begin{aligned}
        \left|
        M\bar\rho_{Q_M}|R|-M\int_R\rho
        \right|
        &\le M\int_R |\rho-\bar\rho_{Q_M}| \\
        &\le C C_0 M r |R| \\
        &\le C C_0 M r^3.
\end{aligned}
\]
By the choice of $r$, this is $C C_0A^3\log M$.  Choose $A$ sufficiently small that $C C_0A^3\le c_S/16$.  Combining with \eqref{eq:local-schmidt},
\[
        \Delta\ge \frac{c_S}{4}\log M-\Delta-\frac{c_S}{16}\log M.
\]
Hence
\[
        2\Delta\ge \frac{3c_S}{16}\log M,
\]
which proves the theorem, after adjusting constants.
\end{proof}

\end{document}
